% Appendix B body (include after \section{Evaluation Codebook} with % Appendix B body (include after \section{Evaluation Codebook} with % Appendix B body (include after \section{Evaluation Codebook} with % Appendix B body (include after \section{Evaluation Codebook} with \input{codebook_appendix})
The codebook below was given to the two human annotators (CS educators) and it mirrors the definitions used for the LLM judge.
\paragraph{Annotation protocol.} Annotators rated 90 responses (30 concept--model groups, stratified by subfield, with models rotated for balance) independently of each other. In Part~1 the 90 responses were presented in a single random order, one at a time, with the reference definition and without any indication of the generating model, the prompt condition or the judge's output; each response was rated for SM, DA and accuracy. Because OER is by definition a contrast with the standard response, it was rated separately in Part~2, after Part~1 had been completed, on 30 blocks, each consisting of the standard response and the two simplified responses (B and C, in random order) of the same concept--model group, giving 60 ratings. Disagreements were resolved by discussion to a consensus label.

\subsection*{Your task}
You will rate LLM-generated explanations of computer science concepts. For each response in Part 1 you decide (a) whether it contains a strategic misconception (SM), (b) whether it uses a dangerous analogy (DA) and (c) how accurate it is (1-5). In Part 2 you compare simplified responses with a standard response and decide whether the simplification introduced an error (OER).


\subsection*{General rules}
\begin{itemize}
  \item Judge technical content only. Ignore style, tone, length, formatting, friendliness and emojis.
  \item The reference definition is your anchor, but it is only one formulation. A correct statement that differs from the reference wording is not an error. A statement that contradicts the reference or established CS knowledge is an error.
  \item Do not penalise omissions as such. An omission counts only if it makes the explanation misleading, i.e. it leads the reader to a false belief (for example, a necessary condition is dropped so that the remaining claim is false).
  \item Rate every response on its own. Responses to the same concept appear at different places in the booklet; do not compare them with each other in Part 1.
  \item Decide from the text. Do not use AI tools. You may consult textbooks to verify a technical fact.
  \item Work independently of the other rater and do not discuss items until the adjudication meeting.
\end{itemize}

\subsection*{Strategic misconception (SM): 0 or 1}
SM = 1 if the response contains at least one conceptually significant error that is stated plausibly and coherently, so that a student would not notice it and would build a wrong mental model of the concept (what it is, how it works, when it applies, or how it differs from a neighbouring concept). Otherwise SM = 0.

\begin{itemize}
  \item NOT an SM: typos, a wrong date or name, harmless imprecision, loose wording that is not false, incompleteness that does not mislead.
  \item SM = 1 example: 'A deadlock is when processes simply wait for a long time.' (distorts the structure of resource holding and circular wait).
  \item SM = 1 example: 'Segmentation eliminates external fragmentation.' (false; variable-sized segments cause it).
  \item SM = 0 example: a wrong year for when an algorithm was published (the concept is unaffected).
  \item If you are unsure whether an error is plausible and consequential, choose 1 and add a short comment.
\end{itemize}

\subsection*{Dangerous analogy (DA): 0 or 1}
DA = 1 if the response uses an analogy, metaphor or everyday comparison that is intuitive but misrepresents a key technical property of the concept, i.e. it invites an inference about the concept that is false.

\begin{itemize}
  \item Ask: what would a student infer from this analogy that is false for the concept? If that false inference concerns a key property, DA = 1.
  \item An analogy is not dangerous merely because it is incomplete. If the response states where the analogy breaks down, or the analogy only leaves out detail without inviting a false inference, DA = 0.
  \item No analogy in the response: DA = 0.
  \item DA = 1 example: a semaphore as a traffic light (invites the inference of a simple binary signal with no counter and no atomic operations).
  \item DA and SM are independent: a response may contain a dangerous analogy without any SM, and vice versa.
\end{itemize}

\subsection*{Accuracy score: 1 to 5}
\begin{itemize}
  \item 5 = fully accurate; no errors or misleading statements.
  \item 4 = accurate overall; minor imprecision or a small omission that does not mislead.
  \item 3 = partly accurate; a notable omission or inaccuracy that could mislead.
  \item 2 = major error(s); an important part of the core concept is wrong.
  \item 1 = fundamentally wrong, or does not address the concept.
\end{itemize}

\subsection*{Part 2: Oversimplification error (OER): 0 or 1}
For each block you see the reference definition, one STANDARD response (written for a prompt asking for an accurate explanation for a first-year student) and two SIMPLIFIED responses, labelled a and b (written for prompts asking for simplification for a student with reading or attention difficulties). Rate each simplified response separately (two ratings per block). The standard response is only the reference point.

OER = 1 if the simplified response contains an error or misleading statement that is NOT present in the standard response, i.e. an error that exists only because of the simplification. Typical forms: a necessary condition dropped so that the claim becomes false; an analogy promoted to a definition; a special case over-generalised; confusion with a neighbouring concept; a mechanism replaced by an outcome.

\begin{itemize}
  \item OER = 0 if the simplified response is merely shorter or less detailed but contains no false or misleading statement.
  \item OER = 0 if the same error is already present in the standard response.
\end{itemize}

\subsection*{Procedure}
\begin{itemize}
  \item Complete Part 1 entirely before you open the Part 2 booklet.
  \item Enter your ratings only in the Excel workbook (sheet Part1, then sheet Part2\_OER). Leave no cell empty.
  \item The comment column is optional but very helpful: write one short sentence for every SM = 1 or DA = 1, and whenever you hesitated.
  \item Expected time: roughly 4-6 hours for Part 1 and 1-2 hours for Part 2. You can work in several sessions.
\end{itemize}
)
The codebook below was given to the two human annotators (CS educators) and it mirrors the definitions used for the LLM judge.
\paragraph{Annotation protocol.} Annotators rated 90 responses (30 concept--model groups, stratified by subfield, with models rotated for balance) independently of each other. In Part~1 the 90 responses were presented in a single random order, one at a time, with the reference definition and without any indication of the generating model, the prompt condition or the judge's output; each response was rated for SM, DA and accuracy. Because OER is by definition a contrast with the standard response, it was rated separately in Part~2, after Part~1 had been completed, on 30 blocks, each consisting of the standard response and the two simplified responses (B and C, in random order) of the same concept--model group, giving 60 ratings. Disagreements were resolved by discussion to a consensus label.

\subsection*{Your task}
You will rate LLM-generated explanations of computer science concepts. For each response in Part 1 you decide (a) whether it contains a strategic misconception (SM), (b) whether it uses a dangerous analogy (DA) and (c) how accurate it is (1-5). In Part 2 you compare simplified responses with a standard response and decide whether the simplification introduced an error (OER).


\subsection*{General rules}
\begin{itemize}
  \item Judge technical content only. Ignore style, tone, length, formatting, friendliness and emojis.
  \item The reference definition is your anchor, but it is only one formulation. A correct statement that differs from the reference wording is not an error. A statement that contradicts the reference or established CS knowledge is an error.
  \item Do not penalise omissions as such. An omission counts only if it makes the explanation misleading, i.e. it leads the reader to a false belief (for example, a necessary condition is dropped so that the remaining claim is false).
  \item Rate every response on its own. Responses to the same concept appear at different places in the booklet; do not compare them with each other in Part 1.
  \item Decide from the text. Do not use AI tools. You may consult textbooks to verify a technical fact.
  \item Work independently of the other rater and do not discuss items until the adjudication meeting.
\end{itemize}

\subsection*{Strategic misconception (SM): 0 or 1}
SM = 1 if the response contains at least one conceptually significant error that is stated plausibly and coherently, so that a student would not notice it and would build a wrong mental model of the concept (what it is, how it works, when it applies, or how it differs from a neighbouring concept). Otherwise SM = 0.

\begin{itemize}
  \item NOT an SM: typos, a wrong date or name, harmless imprecision, loose wording that is not false, incompleteness that does not mislead.
  \item SM = 1 example: 'A deadlock is when processes simply wait for a long time.' (distorts the structure of resource holding and circular wait).
  \item SM = 1 example: 'Segmentation eliminates external fragmentation.' (false; variable-sized segments cause it).
  \item SM = 0 example: a wrong year for when an algorithm was published (the concept is unaffected).
  \item If you are unsure whether an error is plausible and consequential, choose 1 and add a short comment.
\end{itemize}

\subsection*{Dangerous analogy (DA): 0 or 1}
DA = 1 if the response uses an analogy, metaphor or everyday comparison that is intuitive but misrepresents a key technical property of the concept, i.e. it invites an inference about the concept that is false.

\begin{itemize}
  \item Ask: what would a student infer from this analogy that is false for the concept? If that false inference concerns a key property, DA = 1.
  \item An analogy is not dangerous merely because it is incomplete. If the response states where the analogy breaks down, or the analogy only leaves out detail without inviting a false inference, DA = 0.
  \item No analogy in the response: DA = 0.
  \item DA = 1 example: a semaphore as a traffic light (invites the inference of a simple binary signal with no counter and no atomic operations).
  \item DA and SM are independent: a response may contain a dangerous analogy without any SM, and vice versa.
\end{itemize}

\subsection*{Accuracy score: 1 to 5}
\begin{itemize}
  \item 5 = fully accurate; no errors or misleading statements.
  \item 4 = accurate overall; minor imprecision or a small omission that does not mislead.
  \item 3 = partly accurate; a notable omission or inaccuracy that could mislead.
  \item 2 = major error(s); an important part of the core concept is wrong.
  \item 1 = fundamentally wrong, or does not address the concept.
\end{itemize}

\subsection*{Part 2: Oversimplification error (OER): 0 or 1}
For each block you see the reference definition, one STANDARD response (written for a prompt asking for an accurate explanation for a first-year student) and two SIMPLIFIED responses, labelled a and b (written for prompts asking for simplification for a student with reading or attention difficulties). Rate each simplified response separately (two ratings per block). The standard response is only the reference point.

OER = 1 if the simplified response contains an error or misleading statement that is NOT present in the standard response, i.e. an error that exists only because of the simplification. Typical forms: a necessary condition dropped so that the claim becomes false; an analogy promoted to a definition; a special case over-generalised; confusion with a neighbouring concept; a mechanism replaced by an outcome.

\begin{itemize}
  \item OER = 0 if the simplified response is merely shorter or less detailed but contains no false or misleading statement.
  \item OER = 0 if the same error is already present in the standard response.
\end{itemize}

\subsection*{Procedure}
\begin{itemize}
  \item Complete Part 1 entirely before you open the Part 2 booklet.
  \item Enter your ratings only in the Excel workbook (sheet Part1, then sheet Part2\_OER). Leave no cell empty.
  \item The comment column is optional but very helpful: write one short sentence for every SM = 1 or DA = 1, and whenever you hesitated.
  \item Expected time: roughly 4-6 hours for Part 1 and 1-2 hours for Part 2. You can work in several sessions.
\end{itemize}
)
The codebook below was given to the two human annotators (CS educators) and it mirrors the definitions used for the LLM judge.
\paragraph{Annotation protocol.} Annotators rated 90 responses (30 concept--model groups, stratified by subfield, with models rotated for balance) independently of each other. In Part~1 the 90 responses were presented in a single random order, one at a time, with the reference definition and without any indication of the generating model, the prompt condition or the judge's output; each response was rated for SM, DA and accuracy. Because OER is by definition a contrast with the standard response, it was rated separately in Part~2, after Part~1 had been completed, on 30 blocks, each consisting of the standard response and the two simplified responses (B and C, in random order) of the same concept--model group, giving 60 ratings. Disagreements were resolved by discussion to a consensus label.

\subsection*{Your task}
You will rate LLM-generated explanations of computer science concepts. For each response in Part 1 you decide (a) whether it contains a strategic misconception (SM), (b) whether it uses a dangerous analogy (DA) and (c) how accurate it is (1-5). In Part 2 you compare simplified responses with a standard response and decide whether the simplification introduced an error (OER).


\subsection*{General rules}
\begin{itemize}
  \item Judge technical content only. Ignore style, tone, length, formatting, friendliness and emojis.
  \item The reference definition is your anchor, but it is only one formulation. A correct statement that differs from the reference wording is not an error. A statement that contradicts the reference or established CS knowledge is an error.
  \item Do not penalise omissions as such. An omission counts only if it makes the explanation misleading, i.e. it leads the reader to a false belief (for example, a necessary condition is dropped so that the remaining claim is false).
  \item Rate every response on its own. Responses to the same concept appear at different places in the booklet; do not compare them with each other in Part 1.
  \item Decide from the text. Do not use AI tools. You may consult textbooks to verify a technical fact.
  \item Work independently of the other rater and do not discuss items until the adjudication meeting.
\end{itemize}

\subsection*{Strategic misconception (SM): 0 or 1}
SM = 1 if the response contains at least one conceptually significant error that is stated plausibly and coherently, so that a student would not notice it and would build a wrong mental model of the concept (what it is, how it works, when it applies, or how it differs from a neighbouring concept). Otherwise SM = 0.

\begin{itemize}
  \item NOT an SM: typos, a wrong date or name, harmless imprecision, loose wording that is not false, incompleteness that does not mislead.
  \item SM = 1 example: 'A deadlock is when processes simply wait for a long time.' (distorts the structure of resource holding and circular wait).
  \item SM = 1 example: 'Segmentation eliminates external fragmentation.' (false; variable-sized segments cause it).
  \item SM = 0 example: a wrong year for when an algorithm was published (the concept is unaffected).
  \item If you are unsure whether an error is plausible and consequential, choose 1 and add a short comment.
\end{itemize}

\subsection*{Dangerous analogy (DA): 0 or 1}
DA = 1 if the response uses an analogy, metaphor or everyday comparison that is intuitive but misrepresents a key technical property of the concept, i.e. it invites an inference about the concept that is false.

\begin{itemize}
  \item Ask: what would a student infer from this analogy that is false for the concept? If that false inference concerns a key property, DA = 1.
  \item An analogy is not dangerous merely because it is incomplete. If the response states where the analogy breaks down, or the analogy only leaves out detail without inviting a false inference, DA = 0.
  \item No analogy in the response: DA = 0.
  \item DA = 1 example: a semaphore as a traffic light (invites the inference of a simple binary signal with no counter and no atomic operations).
  \item DA and SM are independent: a response may contain a dangerous analogy without any SM, and vice versa.
\end{itemize}

\subsection*{Accuracy score: 1 to 5}
\begin{itemize}
  \item 5 = fully accurate; no errors or misleading statements.
  \item 4 = accurate overall; minor imprecision or a small omission that does not mislead.
  \item 3 = partly accurate; a notable omission or inaccuracy that could mislead.
  \item 2 = major error(s); an important part of the core concept is wrong.
  \item 1 = fundamentally wrong, or does not address the concept.
\end{itemize}

\subsection*{Part 2: Oversimplification error (OER): 0 or 1}
For each block you see the reference definition, one STANDARD response (written for a prompt asking for an accurate explanation for a first-year student) and two SIMPLIFIED responses, labelled a and b (written for prompts asking for simplification for a student with reading or attention difficulties). Rate each simplified response separately (two ratings per block). The standard response is only the reference point.

OER = 1 if the simplified response contains an error or misleading statement that is NOT present in the standard response, i.e. an error that exists only because of the simplification. Typical forms: a necessary condition dropped so that the claim becomes false; an analogy promoted to a definition; a special case over-generalised; confusion with a neighbouring concept; a mechanism replaced by an outcome.

\begin{itemize}
  \item OER = 0 if the simplified response is merely shorter or less detailed but contains no false or misleading statement.
  \item OER = 0 if the same error is already present in the standard response.
\end{itemize}

\subsection*{Procedure}
\begin{itemize}
  \item Complete Part 1 entirely before you open the Part 2 booklet.
  \item Enter your ratings only in the Excel workbook (sheet Part1, then sheet Part2\_OER). Leave no cell empty.
  \item The comment column is optional but very helpful: write one short sentence for every SM = 1 or DA = 1, and whenever you hesitated.
  \item Expected time: roughly 4-6 hours for Part 1 and 1-2 hours for Part 2. You can work in several sessions.
\end{itemize}
)
The codebook below was given to the two human annotators (CS educators) and it mirrors the definitions used for the LLM judge.
\paragraph{Annotation protocol.} Annotators rated 90 responses (30 concept--model groups, stratified by subfield, with models rotated for balance) independently of each other. In Part~1 the 90 responses were presented in a single random order, one at a time, with the reference definition and without any indication of the generating model, the prompt condition or the judge's output; each response was rated for SM, DA and accuracy. Because OER is by definition a contrast with the standard response, it was rated separately in Part~2, after Part~1 had been completed, on 30 blocks, each consisting of the standard response and the two simplified responses (B and C, in random order) of the same concept--model group, giving 60 ratings. Disagreements were resolved by discussion to a consensus label.

\subsection*{Your task}
You will rate LLM-generated explanations of computer science concepts. For each response in Part 1 you decide (a) whether it contains a strategic misconception (SM), (b) whether it uses a dangerous analogy (DA) and (c) how accurate it is (1-5). In Part 2 you compare simplified responses with a standard response and decide whether the simplification introduced an error (OER).


\subsection*{General rules}
\begin{itemize}
  \item Judge technical content only. Ignore style, tone, length, formatting, friendliness and emojis.
  \item The reference definition is your anchor, but it is only one formulation. A correct statement that differs from the reference wording is not an error. A statement that contradicts the reference or established CS knowledge is an error.
  \item Do not penalise omissions as such. An omission counts only if it makes the explanation misleading, i.e. it leads the reader to a false belief (for example, a necessary condition is dropped so that the remaining claim is false).
  \item Rate every response on its own. Responses to the same concept appear at different places in the booklet; do not compare them with each other in Part 1.
  \item Decide from the text. Do not use AI tools. You may consult textbooks to verify a technical fact.
  \item Work independently of the other rater and do not discuss items until the adjudication meeting.
\end{itemize}

\subsection*{Strategic misconception (SM): 0 or 1}
SM = 1 if the response contains at least one conceptually significant error that is stated plausibly and coherently, so that a student would not notice it and would build a wrong mental model of the concept (what it is, how it works, when it applies, or how it differs from a neighbouring concept). Otherwise SM = 0.

\begin{itemize}
  \item NOT an SM: typos, a wrong date or name, harmless imprecision, loose wording that is not false, incompleteness that does not mislead.
  \item SM = 1 example: 'A deadlock is when processes simply wait for a long time.' (distorts the structure of resource holding and circular wait).
  \item SM = 1 example: 'Segmentation eliminates external fragmentation.' (false; variable-sized segments cause it).
  \item SM = 0 example: a wrong year for when an algorithm was published (the concept is unaffected).
  \item If you are unsure whether an error is plausible and consequential, choose 1 and add a short comment.
\end{itemize}

\subsection*{Dangerous analogy (DA): 0 or 1}
DA = 1 if the response uses an analogy, metaphor or everyday comparison that is intuitive but misrepresents a key technical property of the concept, i.e. it invites an inference about the concept that is false.

\begin{itemize}
  \item Ask: what would a student infer from this analogy that is false for the concept? If that false inference concerns a key property, DA = 1.
  \item An analogy is not dangerous merely because it is incomplete. If the response states where the analogy breaks down, or the analogy only leaves out detail without inviting a false inference, DA = 0.
  \item No analogy in the response: DA = 0.
  \item DA = 1 example: a semaphore as a traffic light (invites the inference of a simple binary signal with no counter and no atomic operations).
  \item DA and SM are independent: a response may contain a dangerous analogy without any SM, and vice versa.
\end{itemize}

\subsection*{Accuracy score: 1 to 5}
\begin{itemize}
  \item 5 = fully accurate; no errors or misleading statements.
  \item 4 = accurate overall; minor imprecision or a small omission that does not mislead.
  \item 3 = partly accurate; a notable omission or inaccuracy that could mislead.
  \item 2 = major error(s); an important part of the core concept is wrong.
  \item 1 = fundamentally wrong, or does not address the concept.
\end{itemize}

\subsection*{Part 2: Oversimplification error (OER): 0 or 1}
For each block you see the reference definition, one STANDARD response (written for a prompt asking for an accurate explanation for a first-year student) and two SIMPLIFIED responses, labelled a and b (written for prompts asking for simplification for a student with reading or attention difficulties). Rate each simplified response separately (two ratings per block). The standard response is only the reference point.

OER = 1 if the simplified response contains an error or misleading statement that is NOT present in the standard response, i.e. an error that exists only because of the simplification. Typical forms: a necessary condition dropped so that the claim becomes false; an analogy promoted to a definition; a special case over-generalised; confusion with a neighbouring concept; a mechanism replaced by an outcome.

\begin{itemize}
  \item OER = 0 if the simplified response is merely shorter or less detailed but contains no false or misleading statement.
  \item OER = 0 if the same error is already present in the standard response.
\end{itemize}

\subsection*{Procedure}
\begin{itemize}
  \item Complete Part 1 entirely before you open the Part 2 booklet.
  \item Enter your ratings only in the Excel workbook (sheet Part1, then sheet Part2\_OER). Leave no cell empty.
  \item The comment column is optional but very helpful: write one short sentence for every SM = 1 or DA = 1, and whenever you hesitated.
  \item Expected time: roughly 4-6 hours for Part 1 and 1-2 hours for Part 2. You can work in several sessions.
\end{itemize}
